\section{Agent 工程实践：规划、执行与验证}
\subsection{Plan-Act-Verify 细节}
在实际工程中，Agent 必须具备明确的行动循环：先阅读 Issue、调用检索服务锁定上下文、规划补丁、执行代码修改、运行测试、再利用日志反馈调整。Graham 强调多个工具的组合：LLM 侧重生成，脚本负责操作，测试框架验证结果，最后把证据上传到 dashboard 供人工审核。

他指出，在每一轮 Plan 中应该先生成任务列表，显式地列出哪些文件会被修改、哪些测试要跑。而 Act 阶段则绑定具体的脚本、终端命令，让 Agent 实际写码并执行，Verify 阶段运行测试与静态分析，失败的样本会触发重新 Plan。

\subsection{人机协作与工具覆盖}
\begin{figure}[H]
\centering
\includegraphics[width=\textwidth]{frame-03-human-loop.jpg}
\caption{人类监督者可以在 Agent 运行期间暂停、检查决策并干预，借此纠正实时错误。 \protect\footnotemark}
\end{figure}
\footnotetext{视频画面时间区间：00:32:05--00:32:10。}

OpenHands 允许开发者在任何一步插入人工审查：当 Agent 计划执行危险命令（如部署、数据库写入）时，可以暂停流程，查看日志，并在必要时修改提示或手动执行部分命令。长时间运行的任务也会发送 checkpoint，便于追踪状态。

\begin{importantbox}{最佳实践}
为 Agent 实现 pipeline：先 plan（任务分段、找到依赖），再 act（代码编辑、编译、测试），最后 verify（日志、快照、回滚能力）。失败的步骤需要回溯并把关键上下文传给 LLM。
\end{importantbox}

\begin{knowledgebox}{人工干预的价值}
让人类在 Agent 执行时可以观察其行为，并在将要执行危险操作前暂停，是提升可靠性的关键，尤其在发布构建或数据库操作时。
\end{knowledgebox}

\begin{warningbox}{上下文窗口与幻觉}
Agent 的上下文有限，如果一次性送入整个仓库结构会超出可读窗口，必须先做信息抽取／检索，否则 LLM 会因缺少环境信息产生错觉性的 patch。
\end{warningbox}

\subsection{本章小结}
\begin{itemize}
    \item 规划-执行-验证的循环让 Agent 在工程流程中可控。
    \item 人机协作（暂停、查看日志、跳转上下文）是避免重大误操作的必要手段。
\end{itemize}

